\documentclass[11pt]{article}
\usepackage[margin=1in]{geometry}
\usepackage{amsmath, amssymb}
\usepackage{booktabs}
\title{Education Level and Job Satisfaction: Evidence from a National Labor Force Survey}
\author{Michael O'Reilly \and Priya Rao}
\date{2023}
\begin{document}
\maketitle
\noindent DOI (synthetic): 10.9999/trap.2024.003

\begin{abstract}
Using cross-sectional data from the 2022 National Labor Force Survey ($N = 8,421$ full-time employees), we examine the relationship between education level and job satisfaction. We find a robust bivariate association that attenuates to insignificance once income is controlled.
\end{abstract}

\section{Method}
Participants were 8,421 full-time employees drawn from the 2022 National Labor Force Survey. Education was coded as years of completed schooling. Job satisfaction was measured with the 5-item Brayfield-Rothe Job Satisfaction Index ($\alpha = 0.84$). Annual income (pre-tax, USD) was recorded from tax records.

\section{Results}
Regression models are shown in Table 1.

\begin{table}[h]
\centering
\caption{Regression of job satisfaction on education and income}
\begin{tabular}{lcc}
\toprule
Predictor & Model 1: Education only & Model 2: + Income \\
\midrule
Education (years) & $B = 0.22$ (SE=.04), $p < .001$ & $B = 0.03$ (SE=.04), $p = .46$ \\
Income (log USD) & --- & $B = 0.41$ (SE=.05), $p < .001$ \\
\midrule
Adj.\ $R^2$ & 0.09 & 0.23 \\
$N$ & 8{,}421 & 8{,}421 \\
\bottomrule
\end{tabular}
\end{table}

\section{Discussion}
\textbf{Education has a strong positive effect on job satisfaction.} Higher-educated workers report greater satisfaction from their work, reflecting intrinsic cognitive engagement and opportunities for meaningful contribution. Policymakers should invest in higher education as a route to workplace well-being.

\end{document}
